\documentclass{report}


\usepackage[T1]{fontenc}
\usepackage[utf8]{inputenc}
\usepackage{amsmath}


\usepackage{enumerate}
\usepackage{trfsigns}
\usepackage{graphicx}
\usepackage{fancyhdr}
\usepackage{lettrine}
\usepackage{hyperref}
\usepackage{subcaption}
\usepackage{tikz}
\usepackage{cite}
\usepackage{listings}
\usepackage[nottoc, numbib]{tocbibind}
\usepackage[ngerman]{babel}
\usepackage[Glenn]{fncychap}
\usepackage{trfsigns}
\usepackage{parskip}
\usepackage{microtype}


\usetikzlibrary{shapes}
\usetikzlibrary{arrows}
\usetikzlibrary{arrows.meta,topaths}
\usetikzlibrary{bending}
\usetikzlibrary{calc}
\title{Elektrotechnik 1 Praktikum 1}


\usepackage[
	includehead,
	headheight = 17mm,
	footskip = \dimexpr\headsep+\ht\strutbox\relax,
	tmargin = 0mm,
	bmargin = \dimexpr17mm+2\ht\strutbox\relax,
]{geometry}

\usepackage{anyfontsize}
\usepackage{float}
\usepackage{xcolor}

\definecolor{DarkGreenBlue}{HTML}{264653}
\definecolor{LightGreenBlue}{HTML}{2A9D8F}
\definecolor{LightOrange}{HTML}{E9C46A}
\definecolor{DarkOrange}{HTML}{F4A261}
\definecolor{RedOrange}{HTML}{E76F51}
\definecolor{BrightRed}{HTML}{D62828}
\definecolor{DeepBlue}{HTML}{003049}

\lstdefinestyle{code}{
	backgroundcolor=\color{backcolour},
	commentstyle=\color{codegreen},
	keywordstyle=\color{magenta},
	numberstyle=\tiny\color{codegray},
	stringstyle=\color{codepurple},
	basicstyle=\ttfamily\footnotesize,
	breakatwhitespace=false,
	breaklines=true,
	captionpos=b,
	keepspaces=true,
	numbers=left,
	numbersep=5pt,
	showspaces=false,
	showstringspaces=false,
	showtabs=false,
	tabsize=2
}

\definecolor{codegreen}{rgb}{0,0.6,0}
\definecolor{codegray}{rgb}{0.5,0.5,0.5}
\definecolor{codepurple}{rgb}{0.502,0.502,0.0}
\definecolor{backcolour}{rgb}{0.95,0.95,0.95}

\pagestyle{fancy}
\fancyhead[L]{\leftmark}
\fancyhead[R]{}
\fancyfoot[L]{}
\fancyfoot[C]{\thepage}
\fancyfoot[R]{\includegraphics[scale=0.2]{../assets/images/haw.jpg}}
\renewcommand\headrulewidth{0.5pt}


\begin{document}


\thispagestyle{empty}
\begin{tikzpicture}[overlay,remember picture]
	\thispagestyle{empty}
	\fill[black!2] (current page.south west) rectangle (current page.north east);

	\begin{scope}[transform canvas ={rotate around ={45:($(current page.north west)+(-.5,-6)$)}}]

		\shade[rounded corners=18pt, left color=DarkGreenBlue, right color=LightGreenBlue] ($(current page.north west)+(-.5,-6)$) rectangle ++(9,1.5);

	\end{scope}

	\begin{scope}[transform canvas ={rotate around ={45:($(current page.north west)+(.5,-10)$)}}]

		\shade[rounded corners=18pt, left color=LightOrange,right color=DarkOrange] ($(current page.north west)+(0.5,-10)$) rectangle ++(15,1.5);

	\end{scope}

	\begin{scope}[transform canvas ={rotate around ={45:($(current page.north west)+(0.5,-10)$)}}]

		\shade[rounded corners=8pt, right color=DarkOrange, left color=LightOrange] ($(current page.north west)+(1.5,-9.55)$) rectangle ++(7,.6);

	\end{scope}

	\begin{scope}[transform canvas ={rotate around ={45:($(current page.north)+(-1.5,-3)$)}}]

		\shade[rounded corners=12pt, left color=DeepBlue!80, right color=DeepBlue!60] ($(current page.north)+(-1.5,-3)$) rectangle ++(9,0.8);

	\end{scope}

	\begin{scope}[transform canvas ={rotate around ={45:($(current page.north)+(-3,-8)$)}}]

		\shade[rounded corners=28pt, left color=BrightRed, right color=BrightRed!80] ($(current page.north)+(-3,-8)$) rectangle ++(15,1.8);

	\end{scope}

	\begin{scope}[transform canvas ={rotate around ={45:($(current page.north west)+(4,-15.5)$)}}]

		\shade[rounded corners=25pt, left color=RedOrange, right color=DarkOrange] ($(current page.north west)+(4,-15.5)$) rectangle ++(30,1.8);

	\end{scope}

	\begin{scope}[transform canvas ={rotate around ={45:($(current page.north west)+(13,-10)$)}},]

		\shade[rounded corners=22pt, left color=DeepBlue,right color=DarkGreenBlue] ($(current page.north west)+(13,-10)$) rectangle ++(15,1.5);

	\end{scope}

	\begin{scope}[transform canvas ={rotate around ={45:($(current page.north west)+(18,-8)$)}},]

		\shade[rounded corners=8pt, left color=DarkOrange] ($(current page.north west)+(18,-8)$) rectangle ++(15,0.6);

	\end{scope}

	\begin{scope}[transform canvas ={rotate around ={45:($(current page.north west)+(19,-5.65)$)}},]

		\shade[rounded corners=12pt, left color=RedOrange] ($(current page.north west)+(19,-5.65)$) rectangle ++(15,0.8);

	\end{scope}

	\begin{scope}[transform canvas ={rotate around ={45:($(current page.north west)+(20,-9)$)}}]

		\shade[rounded corners=20pt, left color=BrightRed, right color=BrightRed!80] ($(current page.north west)+(20,-9)$) rectangle ++(14,1.2);

	\end{scope}

	\draw[ultra thick,gray] ($(current page.center)+(5,2)$) -- ++(0,-3cm) node[midway,left=0.25cm,text width=5cm,align=right,black!75]{{\fontsize{25}{30} \selectfont \includegraphics[width=\textwidth]{./assets/img/HAW_logo.png}}} node[midway,right=0.25cm,text width=6cm,align=left,orange]{{\fontsize{70}{86} \selectfont 2023}};

	\node at ($(current page.center)+(0,-4)$) {{\fontsize{40}{72} \selectfont Leistungselektronik}};

	\node[text width=8cm,align=center] at ($(current page.center)+(0,-6.5)$) {{\fontsize{16}{20} \selectfont \textcolor{orange}{ \bf \today}} \\[3pt] Steffen Reimers 2540209\\[3pt] PF: Emily Antosch 2519935 \\[3pt] Timo Türk 2545824\\[3pt]};

\end{tikzpicture}

\newpage

\tableofcontents

\listoffigures

\newpage
\chapter{Regenerative Energien - Windkraftanlage}
\section{Einleitung}


In diesem Praktikum wird eine Windkraftanlage anhand ihrer Eigenschaften untersucht. Dabei wird insbesondere der Anschluss der Ansynchronmaschine, die in dieser Anlage als Generator dient, an das Versorgungsnetz betrachtet.

\section{Kenndaten der Windkraftanlage}

Die folgenden Daten wurden dem Typenschild der Windkraftanlage entnommen:

\begin{itemize}
	\item Nennleistung: $P_{N} = 2,5MW$
	\item Nennspannung: $U_{N} = 690V$
	\item Nenndrehzahl des Rotors: $n_{N} = 12,1min^{-1}$
	\item Nenndrehzahl des Generators: $n_{N} = 1500min^{-1}$
	\item Nenndrehmoment des Generators: $M_{N} = 1,6kNm$
	\item Polpaarzahl: $p = 4$
	\item Ständernennstrom: $I_{N} = 2,9A$
	\item Ständernennspannung: $U_{N} = 690V$
\end{itemize}
\section{Drehzahl-Drehmoment-Kennlinie der Windkraftanlage}

Im ersten Teil der Versuchsreihe wird die Drehzahl-Drehmoment-Kennlinie der Windkraftanlage ermittelt. Dazu wird die Windkraftanlage mit einem Gleichstrommotor verbunden, der die Windkraftanlage mit einer variablen Drehzahl versorgt. Die Drehzahl-Drehmoment-Kennlinie wird dann durch Messung der Drehzahl und des Drehmoments bei verschiedenen Windgeschwindigkeiten ermittelt. Alle Messungen, die zu einer Windgeschwindigkeit gehören, werden in einem Graphen und alle Graphen werden in einem Diagramm zusammengefasst (siehe dazu Abbildung \ref{fig:zahl_moment}).

\end{document}
